\section{Conclusiones}

La presente investigación permitió contrastar la dinámica de la informalidad laboral en el estado de Yucatán en dos periodos estructuralmente distintos: el escenario prepandemia del cuarto trimestre de 2018 y el entorno de recuperación consolidada del cuarto trimestre de 2024. A partir de la evidencia estadística y los modelos econométricos estimados, se presentan las siguientes conclusiones:

\begin{description}
    \item[La paradoja de la formalización selectiva:] Se confirma una trayectoria descendente en la tasa de informalidad estatal, la cual transitó del \textbf{65.77\%} al \textbf{60.06\%}. Sin embargo, este descenso no ha sido homogéneo en términos de acceso. El fortalecimiento del coeficiente asociado al ingreso en el modelo Logit (pasando de $-1.03$ a $-1.16$) sugiere que el mercado laboral yucateco se ha vuelto más excluyente: aunque la incidencia general disminuyó, la barrera económica para acceder a la seguridad social es hoy más pronunciada que hace seis años.
    
    \item[El agotamiento del retorno por experiencia:] Un hallazgo disruptivo de este estudio es la nula significancia estadística de la \textbf{edad} como determinante de la formalidad en ambos periodos ($p > 0.05$). Esto indica que, en el contexto laboral actual de la entidad, la acumulación de años de servicio o la experiencia empírica no constituyen mecanismos suficientes para transitar hacia la formalidad. El mercado laboral de Yucatán parece priorizar las credenciales educativas y los niveles de productividad inmediata por encima de la antigüedad cronológica.
    
    \item[La educación superior como escudo estructural:] La estabilidad del coeficiente de la \textbf{instrucción profesional} (superior a $-1.6$ en ambos modelos) posiciona a la educación universitaria como la herramienta de política pública más efectiva para el combate a la precariedad. Los datos demuestran que poseer un título profesional es el determinante más potente para mitigar el riesgo de informalidad, actuando como un factor protector resiliente incluso ante choques macroeconómicos como los observados entre 2018 y 2024.
    
   \item[Implicaciones para la política pública:] Los resultados sugieren que para reducir la informalidad en Yucatán no basta con la generación de empleo \textit{per se}, sino que es imperativo fortalecer la capacidad de pago de los sectores de baja productividad. Dado que el nivel de ingreso es ahora un predictor más rígido de la informalidad, las políticas orientadas a la \textbf{tecnificación sectorial y el valor agregado} tendrán un impacto marginal mayor en la formalización. Solo incrementando la productividad de los micronegocios será posible elevar el piso salarial real, permitiendo que estos absorban los costos de la seguridad social sin quedar excluidos del sector formal por las presiones del salario mínimo.
\end{description}

\vspace{1em}
\noindent En conclusión, el mercado laboral de Yucatán ha experimentado una transición hacia una estructura donde la informalidad está más fuertemente vinculada a la precariedad salarial. Si bien hay una mayor proporción de trabajadores formales, aquellos que permanecen en la informalidad enfrentan barreras de salida más robustas, donde la educación superior se mantiene como el único vehículo garantizado de movilidad hacia el sector formal.

\subsection{Cumplimiento de Objetivos}
Los resultados de la presente investigación permiten confirmar el cumplimiento del objetivo general, al demostrar que los incrementos al salario mínimo han coincidido con una mayor rigidez en el acceso a la formalidad en Yucatán, evidenciada por el aumento en la sensibilidad del modelo ante la variable ingreso. Respecto al objetivo específico, se identifica que la baja escolaridad constituye el principal factor de desplazamiento hacia la informalidad. No obstante, el análisis revela que la vulnerabilidad no está asociada estrictamente a la condición de juventud, sino a la falta de credenciales educativas y capacidad de generación de ingresos, independientemente de la edad del trabajador.
